\begin{conf-abstract}
{Noble gases as artificial tracers for routine hydrogeological investigations}
{Morgan PEEL, Kai SOLANKI, Daniel HUNKELER, Philip BRUNNER, Rolf KIPFER}
{Eawag}
{Noble gases are increasingly adopted as artificial tracers for hydro(geo)logical studies, thanks to advances in portable, high-resolution dissolved gas measurement technologies [1-3]. Their inert, non-toxic, and invisible nature makes them ideal for sensitive environments, such as rivers and drinking water catchment areas, where maintaining water quality and appearance is essential. The application of gases in aqueous environments poses unique challenges compared to traditional tracers, as potential degassing needs to be accounted for, and ideally avoided, during tracer injection. We present a straightforward method for preparing and injecting tracer solutions with accurately controlled dissolved gas concentrations, whilst avoiding tracer loss. A two-step degassing-pure gas saturation approach enables rapid preparation of high-concentration tracer solutions using readily available materials. We applied this method in a riverbank filtration setting to perform well-to-well tracer tests using helium-4 (He-4) and krypton-84 (Kr-84). Known quantities of both tracers were injected into observation wells upstream of a pumping well. A portable mass spectrometer (“miniRUEDI” GE-MIMS [3]) was used for continuous monitoring in the pumping well. Breakthrough curves of He-4 and Kr-84 enabled precise determination of groundwater flow velocities, travel times, and tracer recovery. These findings illustrate how noble gases complement traditional tracer methods in hydrogeological investigations. The method is adaptable to other gases (e.g., Ne, Kr, Xe, light hydrocarbons), significantly expanding the range of tracers available in groundwater management contexts. 1 : Brennwald et al., 2022, Front. Water (4) 2 : Blanc et al., 2024, Wat. Res. (254) 3 : Brennwald et al., 2016, Env. Sci. Techn. (50)}
\end{conf-abstract}
